\chapter{The Machine-Learning Team}\label{ch:ml:the-machine-learning-team}

A model that took a researcher three weeks took the firm five months to trade. Nothing in it was hard. Every handoff between the
people who built it was a queue: the review committee met weekly, the engineers who rewrote the features for production had a
backlog, the validators sent it back twice because the rewrite did not match the research, and paper trading could not be
hurried. This chapter is about the organisation around the models of the previous twenty-seven chapters: who does what, what
passes between them, and how to measure the result. It builds two things. The first is a \omterm{def:ml:the-machine-learning-team:handoff}{handoff specification} that a
continuous-integration job can check, so that the package a researcher hands over either runs in production as it ran in
research or fails a named test. The second is a simulator of the research-to-production pipeline as a network of queues. On the
chapter's pipeline, the median model reaches production 20 weeks after research ends, and 65\% of models are lost on the way,
most of them at the first review; replacing the rewrite by a validated package cuts the median to 11 weeks, one of which is the
researcher's time spent building the package.

\section{Roles: researcher, engineer, platform}

\begin{definition}[Machine-learning engineer, model owner]\label{def:ml:the-machine-learning-team:roles}
A \emph{machine-learning engineer}\index{machine-learning engineer} builds and operates the systems that train, serve and monitor
models in production, combining the skills of data scientists, data engineers and software engineers. The \emph{model
owner}\index{model owner} is the named person accountable for a model in production: its performance, its monitoring, its
retraining and its retirement, and the answers to the questions validators, risk managers and regulators ask about it.
\end{definition}

A trading firm's machine-learning work involves at least three kinds of people. Researchers (quantitative researchers or data
scientists) find the signal and fit the model, with the methods of this book and the craft of Book~7. \omterm{def:ml:the-machine-learning-team:roles}{Machine-learning
engineers} turn the model into a service that runs every day without its author: features computed online as they were offline
(chapter~24), a compiled model inside the latency budget (chapter~26), monitors (chapter~27). A platform team builds what every
model shares. Kreuzberger, Kühl and Hirschl, from a literature review, a tool review and expert interviews, list seven roles
around machine-learning operations (business stakeholder, solution architect, data scientist, data engineer, software engineer,
DevOps engineer, and the machine-learning or MLOps engineer who combines several of them); a trading firm adds the model
validator of Book~6 (chapter~26), independent of the others, and the trader or portfolio manager who uses the model.

The \omterm{def:ml:the-machine-learning-team:roles}{model owner} is a role, not a team. It is usually the researcher who proposed the model or the manager of the desk that uses
it; what matters is that it is one named person, recorded in the registry (chapter~25) and on the \omterm{def:ml:interpretability-and-model-governance:card}{model card} (chapter~21), who
answers the page of a \omterm{def:ml:monitoring-and-retraining:performance-alarm}{performance alarm} and signs the decision to retrain or retire. A model without an owner is the one whose
monitors page nobody. Amershi and co-authors, observing teams at Microsoft, found that model customisation and reuse require
skills different from those of software teams, and that models are harder to treat as distinct modules than ordinary software
components; both are reasons the division of work between researchers and engineers is a design choice with a cost, not a given.

\section{How work moves: the handoff}

\begin{definition}[Handoff specification]\label{def:ml:the-machine-learning-team:handoff}
A \emph{handoff specification}\index{handoff specification} is the contract between the people who build a model and the people
who run it: the list of what a model package must contain (the artefact and its content hash, the \omterm{def:ml:data-and-feature-stores:store}{feature definitions}, test
vectors of inputs with the outputs the model must reproduce, the latency budget with a measurement, the monitoring specification,
the \omterm{def:ml:interpretability-and-model-governance:card}{model card}) and the checks the package must pass before production accepts it.
\end{definition}

In the hook's firm the handoff was a rewrite. The researcher handed over a notebook; an engineer reimplemented the features and
the model in production code; the validator compared the two and, when they differed, sent the model back. Every difference was
a round trip through two queues. Zinkevich's rules for machine-learning engineering say it directly: ``Re-use code between your
training pipeline and your serving pipeline whenever possible''. The \omterm{def:ml:the-machine-learning-team:handoff}{handoff specification} is the other half of that advice: when
the same code runs in both places, what passes between the teams is a package that can be checked mechanically.

The chapter's package is for chapter~26's compiled forest (\cref{lst:ml:team:validate}). Its manifest names the forest file and
its content hash (Book~7's \texttt{firm.workflow}); sixteen \omterm{def:ml:data-and-feature-stores:store}{feature definitions} in \texttt{firm.featstore}'s format; research's
values of those features on a reference session of the simulated order book; 200 test vectors with the forest's outputs, to be
reproduced to $10^{-12}$; a latency budget of 5 microseconds with the measured 2.6; four monitors with thresholds, a false-page
rate and an owner each; and the \omterm{def:ml:interpretability-and-model-governance:card}{model card}'s fifteen fields. The validator runs nine checks, and the clean package passes all of
them. Then the chapter plants eight defects a handoff typically carries, one at a time (\cref{tab:ml:team:defects}).

\begin{table}[htbp]
\centering\small
\begin{tabular}{@{}p{5.4cm}p{6.6cm}@{}}
\toprule
defect planted & checks that fail\\
\midrule
none (the package as handed over) & none of the nine\\
artefact rebuilt after it was hashed & artefact hash; test vectors\\
a feature reads an input the feed does not carry & feature specification\\
a 5-second window written as 5\,000 (milliseconds) & feature vectors\\
one feature missing (15 for a model reading 16) & feature specification\\
test vectors holding another model's outputs & test vectors\\
latency of the library call, not the compiled model & latency\\
a monitor without an owner & monitoring\\
a \omterm{def:ml:interpretability-and-model-governance:card}{model card} without limitations & \omterm{def:ml:interpretability-and-model-governance:card}{model card}\\
\bottomrule
\end{tabular}
\caption{The handoff validator on the clean package and on eight planted defects. Data: \texttt{ml\_team.handoff\_table}.}
\label{tab:ml:team:defects}
\end{table}

Every defect is caught, and by the check written for it, with one lesson in the third row. A window in the wrong unit is still a
valid \omterm{def:ml:data-and-feature-stores:store}{feature definition}, and the offline and online computations agree with each other, since both use the same wrong window:
the specification and parity checks pass. Only the comparison with research's own feature values on a reference session finds
it. Test vectors must cover the features as well as the model.

\section{What the platform owns}

\begin{definition}[Machine-learning platform]\label{def:ml:the-machine-learning-team:platform}
A \emph{machine-learning platform}\index{machine-learning platform} is the shared infrastructure on which every model is built
and run: data and \omterm{def:ml:data-and-feature-stores:store}{feature stores}, training infrastructure, experiment tracking and the \omterm{def:ml:experiment-tracking-and-reproducibility:registry}{model registry}, compilation and serving,
monitoring, and the checks a model package passes on its way to production, owned by one team and used by all.
\end{definition}

The platform owns everything that is the same for every model. In this book's terms: the \omterm{def:ml:data-and-feature-stores:store}{feature store} with its parity tests
(chapter~24), the training infrastructure (chapter~23), the tracker, registry and \omterm{def:ml:experiment-tracking-and-reproducibility:registry}{audit trail} (chapter~25), the compilers and
serving kernels (chapter~26), the monitors (chapter~27) and the handoff validator. Researchers own the model and its features;
the platform owns the guarantee that a feature means the same thing in research and in production. The line matters because each
side tends to push work across it: researchers ask for one-off pipelines, and platforms grow features no model uses. A useful
test is Zinkevich's fourth rule, ``Keep the first model simple and get the infrastructure right'': a platform is justified by the
second model, not the first. Its other justification is visibility: among the risks specific to machine-learning systems,
Sculley and co-authors list undeclared consumers and data dependencies, which a shared registry and \omterm{def:ml:data-and-feature-stores:store}{feature store} make explicit.

The validator is the platform's most leveraged piece, because it turns a conversation into a test. In the chapter's pipeline
below, it is what lets the integration of a model take a week instead of a rewrite of four.

\section{Measuring the pipeline}

The pipeline is a network of queues (\cref{lst:ml:team:sim}); its parameters are illustrative, not measured at any firm. Models arrive when research ends, at 0.35 a week (18 a year), and
pass five stations. The review committee (one slot, 0.6 weeks a model) kills 35\% of the models it reviews and sends 15\% back
for two weeks of research fixes. Two engineers reimplement each surviving model in about four weeks; 30\% of rewrites go back to
the researcher for a week and a half of clarification and are then rewritten, and 5\% are abandoned. One validator (1.5 weeks)
kills 10\% and returns 25\% to the engineers. Paper trading lasts four weeks and kills 20\%; a two-week canary kills 5\%. Service
times are exponential except for the two fixed trading periods; the first 500 of 4\,000 simulated models are discarded while the
pipeline fills.

The simulator is checked on the one case with a closed form. A single station of the M/M/1 queue of Book~4 (chapter~8), arrivals
at 0.8 and service at 1 a week, gives a mean time in the system of 5.05 weeks against the theoretical $1/(\mu-\lambda)=5$; the
stationary law of the corresponding birth--death chain, computed with Book~4's \texttt{firm.queues}, gives a mean of 4.00 models in
the system and an empty system 20\% of the time, and the simulation's arrival rate times its mean time gives 4.04, as Little's law
$L=\lambda W$ requires. On the whole pipeline, the time-average number of models in it is 4.87 and the arrival rate times their
mean time in it (14.0 weeks, counting killed models until they die) is 4.90.

\begin{table}[htbp]
\centering\small
\begin{tabular}{@{}lcc@{}}
\toprule
 & rewrite (handoff) & validated package\\
\midrule
median lead time to production (weeks) & 20.4 & 11.1\\
quartiles (weeks) & 14.8 to 30.2 & 9.6 to 13.7\\
90th percentile (weeks) & 42.9 & 16.5\\
share of models reaching production & 35.1\% & 37.5\%\\
lost at review & 41.1\% & 39.5\%\\
lost at the build (reimplementation / integration) & 5.3\% & 3.5\%\\
lost at validation & 7.9\% & 6.9\%\\
lost in paper trading & 8.7\% & 10.3\%\\
lost in the canary & 1.9\% & 2.2\%\\
passes through the build, per model past review & 1.79 & 1.14\\
engineers' utilisation & 74\% & 12\%\\
\bottomrule
\end{tabular}
\caption{The research-to-production pipeline with a rewrite at the handoff and with a validated package (lead time includes one
week of research spent building the package); 3\,500 models each. Data: \texttt{ml\_team.pipeline}.}\label{tab:ml:team:pipeline}
\end{table}

\Cref{tab:ml:team:pipeline} is the chapter's result. With the rewrite, the median model reaches production 20.4 weeks after
research ends, about five months, and one in ten takes more than 42.9 weeks; 35\% of models get there. The largest loss is at
review, 41\% of all models: that is where most models should die, cheaply and early. Of the time, service without any waiting
would account for 14.9 weeks of the median; the rest is queueing, concentrated at the engineers, who are busy 74\% of the time
because each model passes through them 1.79 times on average. Replacing the rewrite by a validated package (integration in about
a week, rework at the build falling from 30\% to 5\% and at validation from 25\% to 10\%, plus a week of the researcher's time)
cuts the median to 11.1 weeks, a saving of 9.3 weeks (45\%), and the 90th percentile from 42.9 to 16.5. The shares lost at each
gate barely move, which is right: removing a handoff should make decisions faster, not different.

\begin{omfigure}
\begin{tikzpicture}
\begin{axis}[width=11cm,height=5.8cm,xlabel={models arriving per week},ylabel={median lead time (weeks)},xmin=0.19,xmax=0.45,
  ymin=0,ymax=72,tick label style={font=\small},label style={font=\small},grid=major,grid style={gray!25},legend pos=north west,
  legend cell align=left,legend style={font=\small},xticklabel style={/pgf/number format/fixed,/pgf/number format/precision=2},
  table/col sep=comma]
\addplot[omProp,very thick,mark=*,mark size=1.5pt] table[x=rate,y=handoff]{figdata/ml/28-the-machine-learning-team/rate.csv};
\addplot[omDef,very thick,dashed,mark=square*,mark size=1.5pt,mark options={solid}] table[x=rate,y=package]{figdata/ml/28-the-machine-learning-team/rate.csv};
\legend{rewrite at the handoff,validated package}
\end{axis}
\end{tikzpicture}

\omcaption{Median lead time to production against the rate at which research sends models, with a rewrite at the handoff and
with a validated package, means over three simulated pipelines of 4\,000 models. Data: \texttt{ml\_team.rate\_curve}.}
\label{fig:ml:team:rate}
\end{omfigure}

The table holds at one arrival rate; \cref{fig:ml:team:rate} shows why it matters more than it seems. The engineers' station is
a queue near saturation, and a queue's waiting time grows without bound as its utilisation approaches one: from 0.35 to 0.44
models a week, a quarter more research, the median lead time with the rewrite rises from about 22 to 70 weeks, while with the
package it stays near 11. A pipeline is measured by its lead time, the share lost at each gate and the utilisation of its busiest
station, and a manager who adds researchers without looking at the last number makes every model slower.

\begin{method}[Running a research-to-production pipeline]\label{met:ml:team:method}
\begin{enumerate}
\item Name a \omterm{def:ml:the-machine-learning-team:roles}{model owner} for every model, in the registry and on the \omterm{def:ml:interpretability-and-model-governance:card}{model card}.
\item Replace rewrites by packages: the same feature and model code in research and production, and a \omterm{def:ml:the-machine-learning-team:handoff}{handoff specification}
checked in continuous integration, with test vectors for the features as well as the model.
\item Kill models early: the cheapest gate should reject the most.
\item Measure lead time, losses per gate and the utilisation of each team; keep the busiest well below saturation.
\item Put what every model shares on a platform owned by one team, and justify it by the second model.
\end{enumerate}
\end{method}

\section{Tutorial: who owns the model?}\label{tut:ml:the-machine-learning-team:tut}

\begin{tutorial}
\textbf{Goal.} Validate a handoff package in the way continuous integration would, plant defects and see which check catches each;
simulate the pipeline with and without the rewrite. \textbf{End state:} \cref{tab:ml:team:defects}, \cref{tab:ml:team:pipeline},
\cref{fig:ml:team:rate}.
\begin{tutsteps}
\item \textbf{The validator.}
\omcode{code/firm/modelpkg/firm_modelpkg.py}{60}{92}{Schema, artefact hash, feature specification, parity and feature vectors.}\label{lst:ml:team:validate}
\item \textbf{The rest of the checks.}
\omcode{code/firm/modelpkg/firm_modelpkg.py}{93}{105}{Test vectors through the forest kernel, latency, monitoring and the model card.}\label{lst:ml:team:checks}
\item \textbf{The planted defects} (\texttt{ml\_team.\_defects}).
\omcode{code/ml/28-the-machine-learning-team/python/ml_team.py}{96}{119}{Seven of the eight defects, each an edit of a loaded package.}\label{lst:ml:team:defects}
\item \textbf{The pipeline.}
\omcode{code/ml/28-the-machine-learning-team/python/ml_team.py}{145}{154}{Stations, with and without the rewrite.}\label{lst:ml:team:stations}
\item \textbf{The simulator.}
\omcode{code/firm/modelpkg/firm_modelpkg.py}{143}{171}{The event loop: queues, service, and the kill, rework and pass decisions.}\label{lst:ml:team:sim}
\item \textbf{Run} \texttt{ml\_team.handoff\_table()}, \texttt{pipeline(True)}, \texttt{pipeline(False, 1.0)},
\texttt{mm1\_check()} and \texttt{fig\_team.py} (about 20 seconds on one core).
\end{tutsteps}
\textbf{What to change next.} Let the review committee meet twice as often, and find which station is then the busiest.
\end{tutorial}

\section{Build: the model package}\label{bld:ml:the-machine-learning-team:modelpkg}

\begin{build}
\bfield{Purpose} A contract between research and production that a machine can check, and a model of the pipeline it flows through.
\bfield{Interface} \texttt{REQUIRED}, \texttt{INPUTS}, \texttt{MONITORS}; \texttt{load(path)}; \texttt{validate(pkg, events,
times)} returning (check, passed, detail) triples, on \texttt{firm.workflow}, \texttt{firm.featstore}, \texttt{firm.mlinfer} and
\texttt{firm.modelcard}; \texttt{Station(name, servers, mean, kill, rework, rework\_to, dist, next\_to)};
\texttt{simulate\_pipeline(stations, rate, n, seed)}.
\bfield{Rules} Every field of the specification is required; a check that cannot run fails; feature values, not only model
outputs, are compared with research's; the simulator is deterministic given its seed.
\bfield{Acceptance tests} \texttt{code/firm/modelpkg/tests/}: a complete package passes; a missing field, a wrong hash, an
unknown input and a stale test vector each fail their check; one M/M/1 station matches $1/(\mu-\lambda)$; a pure delay line
returns its fixed time; kill and rework probabilities are respected.
\bfield{Stretch} Run the validator as a pre-merge job on the registry, with the package's manifest as the registered artefact;
priorities at the busiest station.
\end{build}

\begin{omsources}
\item S.~Amershi and co-authors, ``Software engineering for machine learning: a case study'', \emph{ICSE Software Engineering
in Practice}, 2019.
\item D.~Kreuzberger, N.~Kühl and S.~Hirschl, ``Machine learning operations (MLOps): overview, definition, and
architecture'', \emph{IEEE Access}, 2023.
\item M.~Zinkevich, \emph{Rules of Machine Learning: Best Practices for ML Engineering}, Google for Developers.
\item J.~D.~C.~Little, ``A proof for the queuing formula $L=\lambda W$'', \emph{Operations Research}, 1961.
\item D.~Sculley and co-authors, ``Hidden technical debt in machine learning systems'', \emph{NeurIPS}, 2015.
\end{omsources}

\section{Exercises}

\begin{exercise}[$\star$]\label{exo:ml:the-machine-learning-team:1}
Name the three kinds of people in a trading firm's machine-learning work and one thing each owns.
\end{exercise}

\begin{exercise}[$\star$]\label{exo:ml:the-machine-learning-team:2}
Why must the \omterm{def:ml:the-machine-learning-team:roles}{model owner} be one named person? What goes wrong when it is a team?
\end{exercise}

\begin{exercise}[$\star$]\label{exo:ml:the-machine-learning-team:3}
Check Little's law on the chapter's single station: arrivals at 0.8 a week, a mean time in the system of 5.05 weeks. What is
the mean number of models in the system?
\end{exercise}

\begin{exercise}[$\star\star$]\label{exo:ml:the-machine-learning-team:4}
Why do the specification and parity checks pass a window written in milliseconds, and what check catches it?
\end{exercise}

\begin{exercise}[$\star\star$]\label{exo:ml:the-machine-learning-team:5}
The engineers are busy 74\% of the time with the rewrite. Using the M/M/1 formula as a rough guide, by how much does the mean
wait change if their utilisation rises to 90\%?
\end{exercise}

\begin{exercise}[$\star\star$]\label{exo:ml:the-machine-learning-team:6}
\emph{Find the flaw.} ``Our pipeline is slow because validation kills too many models; we should relax the validators.''
\end{exercise}

\begin{exercise}[$\star\star\star$]\label{exo:ml:the-machine-learning-team:7}
\emph{Coding.} Keep the rewrite and give the engineers a third member. Compare the median lead time with the package's, and
say which change you would make first.
\end{exercise}

\begin{exercise}[$\star\star\star$]\label{exo:ml:the-machine-learning-team:8}
Write the \omterm{def:ml:the-machine-learning-team:handoff}{handoff specification} for chapter~27's monitored linear model: the fields, the test vectors, and the checks.
\end{exercise}

\section{Problem: Who Owns the Model?}

\begin{problem}[{Weekend problem --- who owns the model?}]\label{pb:ml:the-machine-learning-team:1}
The chapter's handoff package and pipeline.

\medskip\noindent\textbf{Part I --- Roles.}
\begin{enumerate}
\item Name the roles in a trading firm's machine-learning work.
\item What does the \omterm{def:ml:the-machine-learning-team:roles}{model owner} answer for?
\item What does the platform own, and what does it not?
\item Why is the division between researchers and engineers a design choice?
\end{enumerate}

\medskip\noindent\textbf{Part II --- The handoff.}
\begin{enumerate}[resume]
\item What does the chapter's package contain?
\item Which checks does the validator run?
\item Which check catches each planted defect?
\item Why must test vectors cover the features?
\end{enumerate}

\medskip\noindent\textbf{Part III --- The pipeline.}
\begin{enumerate}[resume]
\item Describe the simulated pipeline.
\item How is the simulator checked?
\item Where do models die, and where does the time go?
\item What happens as more models arrive?
\end{enumerate}

\medskip\noindent\textbf{Part IV --- The verdict.}
\begin{enumerate}[resume]
\item State the \emph{named result}: the pipeline's median lead time, the share of models lost at each gate, and the saving from
removing one handoff.
\item Why should the losses per gate barely change when the handoff is removed?
\item What does the researcher pay for the package, and is it worth it?
\item Would a third engineer do as well?
\item What should the firm measure every month?
\item Who should be the owner of the chapter's forest?
\item What should happen when a package fails a check?
\item In one sentence: what makes a handoff cheap?
\end{enumerate}
\end{problem}

\section{Interview questions}

\begin{interviewq}[$\star$ \iqroles{mle}]\label{iq:ml:the-machine-learning-team:1}
What does a \omterm{def:ml:the-machine-learning-team:roles}{machine-learning engineer} do that a researcher does not?
\end{interviewq}

\begin{interviewq}[$\star\star$ \iqroles{mle, developer}]\label{iq:ml:the-machine-learning-team:2}
What should a research team hand to production with a model?
\end{interviewq}

\begin{interviewq}[$\star\star$ \iqroles{mle}]\label{iq:ml:the-machine-learning-team:3}
Your features differ between research and production. How do you find out, and how do you stop it happening again?
\end{interviewq}

\begin{interviewq}[$\star\star$ \iqroles{researcher, mle}]\label{iq:ml:the-machine-learning-team:4}
Models take five months to reach production. How would you find out why?
\end{interviewq}

\begin{interviewq}[$\star\star$ \iqroles{mle}]\label{iq:ml:the-machine-learning-team:5}
What belongs on a \omterm{def:ml:the-machine-learning-team:platform}{machine-learning platform}, and what should stay with the model's team?
\end{interviewq}

\begin{interviewq}[$\star\star\star$ \iqroles{mle}]\label{iq:ml:the-machine-learning-team:6}
Your platform team is asked to double the number of models it supports next year. What do you measure first?
\end{interviewq}
